% Pandoc template for the fed_fraud documents.
% Based on proposal_gen/templates/detailed_proposal_template/detailed_1_classic.tex
% (Classic report): centered title page, table of contents, ruled section
% headings and the Geodesic house footer. Fonts are Latin Modern (lmodern).
% Two layouts: the full report (title page and contents) and, with
% -V brief, a compact one-page-brief header in the same style.
\documentclass[letterpaper,11pt]{article}
\usepackage[T1]{fontenc}
\usepackage[utf8]{inputenc}
\usepackage{lmodern}
\usepackage[letterpaper,top=0.9in,bottom=1.2in,left=1in,right=1in,footskip=0.6in]{geometry}
\usepackage[table]{xcolor}
\usepackage{geodesicbrand}
\usepackage{amsmath,amssymb}
\usepackage{enumitem}
\usepackage{array}
\usepackage{calc}
\usepackage{longtable}
\usepackage{makecell}
\usepackage[most]{tcolorbox}
\usepackage{titlesec}
\usepackage{needspace}
\usepackage{etoolbox}
\usepackage{hyperref}
\hypersetup{colorlinks=true,urlcolor=coolblack,linkcolor=coolblack,filecolor=geoblue,
  pdftitle={Social Security Administration},pdfauthor={\geofirm}}
\urlstyle{same}

\setlist{itemsep=2pt,leftmargin=1.4em}
\setlength{\parindent}{0pt}
\setlength{\parskip}{6pt}
\providecommand{\tightlist}{\setlength{\itemsep}{0pt}\setlength{\parskip}{0pt}}
\setcounter{secnumdepth}{0}

% Section style: house rule under each heading
\titleformat{\section}{\Large\bfseries\color{coolblack}}{\thesection}{0.6em}{}[{\color{coolblack}\titlerule[0.5pt]}]
\titleformat{\subsection}{\large\bfseries\color{geomid}}{\thesubsection}{0.6em}{}
\titleformat{\subsubsection}{\normalsize\bfseries\color{geomid}}{\thesubsubsection}{0.6em}{}
\titlespacing*{\section}{0pt}{10pt}{6pt}

% Tables are written by latex/filters.lua (navy header, banded rows)
\setlength{\LTpre}{6pt}\setlength{\LTpost}{6pt}

\begin{document}

% BRIEF HEADER
\thispagestyle{geodesic}
\begin{minipage}[c]{0.13\textwidth}\includegraphics[width=\linewidth]{icon.png}\end{minipage}\hfill
\begin{minipage}[c]{0.84\textwidth}
  {\textsc{\large \geofirm}}\par\vspace{2pt}
  {\small\color{geogray}Federal fraud expansion\enspace\textperiodcentered\enspace When: \compiletime}
\end{minipage}\par\vspace{6pt}
{\color{coolblack}\rule{\linewidth}{0.3mm}}\par\vspace{2pt}
{\LARGE\bfseries\color{coolblack} Social Security Administration}\par
\vspace{2pt}{\large\color{geomid} One-page brief}\par
\vspace{-2pt}{\color{coolblack}\rule{\linewidth}{0.3mm}}\par\vspace{4pt}
\pagestyle{geodesic}

\textbf{Priority:} 22 of 33 (score 2.7). \textbf{FY2026 outlays:}
\$1,664B (USAspending). \textbf{As of:} 2026-10-27. Public data only.

\hypertarget{bottom-line}{%
\section{Bottom line}\label{bottom-line}}

SSA pays more benefits than any agency but publishes almost nothing at
the payee level. Its error rates are low for retirement benefits and
high for Supplemental Security Income (11.6\%). The most useful
public-data role for SSA is as a supplier: its death records are what
other agencies (OPM, PBGC, CMS) are missing when they pay dead people.
Approach SSA as a data partner first and a target second.

\hypertarget{exposure}{%
\section{Exposure}\label{exposure}}

\begin{itemize}
\tightlist
\item
  SSI: \$7.3B improper (11.6\%). Disability Insurance: \$2.3B. (GAO,
  FY2025)
\item
  PBGC credited a pension plan for 3,479 deceased participants because
  it did not check SSA's death file, overpaying about \$127M (2023).
\item
  GAO found FEHB paying about 400 providers for services dated at least
  50 days after their death (2026).
\end{itemize}

\hypertarget{public-data-we-can-use}{%
\section{Public data we can use}\label{public-data-we-can-use}}

\begin{itemize}
\tightlist
\item
  \textbf{SSA State and county statistics:} beneficiaries and payments
  by program and county; SSI recipients by state and county.
\item
  \textbf{SSA OIG audits and investigations:} representative payee and
  disability fraud cases.
\item
  \textbf{Disability fraud cases} from DOJ and the Cooperative
  Disability Investigations units, as labeled examples.
\end{itemize}

\hypertarget{detection-plays}{%
\section{Detection plays}\label{detection-plays}}

\begin{enumerate}
\def\labelenumi{\arabic{enumi}.}
\tightlist
\item
  \textbf{Geographic outliers:} counties where SSI or disability rolls
  and award rates are far above peers after adjusting for age and
  health.
\item
  \textbf{Facilitator networks:} attorneys, doctors and representative
  payees who recur across charged disability cases.
\item
  \textbf{Death data gaps (cross-agency):} quantify how much other
  agencies' payments could be stopped with full death data.
\end{enumerate}

\hypertarget{first-prototype}{%
\section{First prototype}\label{first-prototype}}

A county-level SSI and disability outlier map, paired with a one-page
case for SSA death data sharing built from the OPM and PBGC findings.

\hypertarget{limits}{%
\section{Limits}\label{limits}}

No public payee-level data, and full death data sharing is restricted by
law. This brief is mainly a data-access proposal.

\hypertarget{who-to-take-it-to}{%
\section{Who to take it to}\label{who-to-take-it-to}}

SSA Office of Anti-Fraud Programs, SSA OIG, and Treasury's Do Not Pay.

\hypertarget{sources}{%
\section{Sources}\label{sources}}

\begin{itemize}
\tightlist
\item
  \href{https://www.gao.gov/products/gao-26-108694}{GAO-26-108694:
  FY2025 improper payments}
\item
  \href{https://www.plansponsor.com/pbgc-says-central-states-will-reimburse-sfa-overpayment/}{PlanSponsor:
  PBGC overpaid Central States \$127M for deceased participants}
\item
  \href{https://files.gao.gov/reports/GAO-26-108139/index.html}{GAO-26-108139:
  FEHB paid deceased and excluded providers}
\item
  \href{https://www.usaspending.gov/agency}{USAspending agency profiles
  (FY2026 outlays)}
\end{itemize}

\end{document}
